%% Section 6: Analysis and Discussion
%% Japanese draft (draft2/06_analysis.md) embedded as comments.
%% English translation written below each paragraph.
%% -----------------------------------------------------------------------

\section{Analysis and Discussion}
\label{sec:analysis}

\subsection{Ablation Study: Contribution of the Modulo Stream}
\label{sec:ablation}

% Modulo ストリームと FiLM 融合の効果を単独で定量化するため、全 3 スケールにわたって IntSeqBERT と
% Magnitude-only のアブレーションモデルを比較する。
% 表4 に Large 設定における差分をまとめる。
To isolate the contribution of the Modulo stream and FiLM fusion, we compare IntSeqBERT with
the Magnitude-only Ablation model across all three scales. Table~\ref{tab:ablation} summarises
the differences at the Large scale.

\begin{table}[t]
\caption{IntSeqBERT vs.\ Ablation at the Large scale (test performance).
         $\Delta$ = IntSeq $-$ Ablation.}
\label{tab:ablation}
\centering
\begin{tabular}{lrrr}
\hline
Metric              & IntSeq & Ablation & $\Delta$ \\
\hline
Mag Acc (\%)        &  95.85 &    89.70 &   $+6.15$ \\
Sign accuracy (\%)  &  98.54 &    98.29 &   $+0.25$ \\
MMA (\%)            &  50.38 &    35.22 &  $+15.16$ \\
mod-2 accuracy (\%) &  85.65 &    72.13 &  $+13.52$ \\
Solver Top-1 (\%)   &  19.09 &    11.75 &   $+7.34$ \\
Magnitude MSE       &  0.142 &     0.371&   $-0.228$\\
\hline
\end{tabular}
\end{table}

% 設計通り、Modulo ストリームは Modulo 予測において最大の改善をもたらす（MMA +15.2pt、パリティ
% +13.5pt）。Magnitude 予測への寄与は比較的小さいが（Acc_{0.5} +6.2pt、MSE -0.228）、一貫して
% 有意である。これは周期的算術構造がスケール回帰に対して相補的な帰納バイアスを提供していることを示す。
% すなわち、x_i の 2 から 101 に対する剰余を知ることは、とりうる Magnitude の範囲を大幅に絞り込む。
% 符号精度の差は小さく（+0.25pt）、符号は Magnitude のみからおおむね推定可能であることを示す。
% Solver 精度は Modulo ストリームによって +7.3pt 改善される。Dense および Sieve モードにおいてより
% 正確な剰余情報がスコアリングを精緻化することが主な要因である。
As designed, the Modulo stream delivers the largest improvement in Modulo prediction (MMA
$+15.2$\,pt, parity $+13.5$\,pt). Its contribution to Magnitude prediction is smaller but
consistently significant (Acc$_{0.5}$ $+6.2$\,pt, MSE $-0.228$), suggesting that periodic
arithmetic structure provides a complementary inductive bias for scale regression: knowing
$x_i \bmod m$ for $m \in \{2, \ldots, 101\}$ substantially constrains the set of plausible
Magnitude values. The sign accuracy gap is small ($+0.25$\,pt), indicating that sign is largely
inferrable from magnitude alone. Solver accuracy improves by $+7.3$\,pt, primarily because
more accurate residue information refines scoring in Dense and Sieve modes.

\subsection{Scaling Behaviour}
\label{sec:scaling}

% 表5 に IntSeqBERT のモデルサイズ別の指標改善幅を示す。
% Magnitude 精度はモデルサイズに対して比較的緩やかな改善に留まる（Small→Large で +1.1pt）のに対し、
% Modulo 精度はより大きな改善を示す（+10.0pt）。
% これは剰余算術がより合成的であり、表現容量の増加からより多くの恩恵を受けやすいという直観と整合する。
% Solver Top-1 精度も Modulo 精度の傾向に追随して一貫した改善を示す（+5.0pt）。
Table~\ref{tab:scale_trend} shows how IntSeqBERT's metrics improve with model size.
Magnitude accuracy improves only modestly from Small to Large ($+1.1$\,pt), whereas Modulo
accuracy improves substantially ($+10.0$\,pt). This is consistent with the intuition that
modular arithmetic is more compositional and thus benefits more from increased representational
capacity. Solver Top-1 accuracy follows the Modulo trend with consistent improvement ($+5.0$\,pt).

\begin{table}[t]
\caption{IntSeqBERT scaling: test-split metrics (Small / Middle / Large).}
\label{tab:scale_trend}
\centering
\begin{tabular}{lrrrr}
\hline
Metric              & Small  & Middle & Large  & $\Delta$ (S$\to$L) \\
\hline
Mag Acc (\%)        &  94.73 &  95.71 &  95.85 & $+1.12$ \\
MMA (\%)            &  40.43 &  46.88 &  50.38 & $+9.95$ \\
mod-2 accuracy (\%) &  81.97 &  84.50 &  85.65 & $+3.68$ \\
Solver Top-1 (\%)   &  14.05 &  17.02 &  19.09 & $+5.04$ \\
Magnitude MSE       &  0.228 &  0.164 &  0.142 & $-0.086$\\
\hline
\end{tabular}
\end{table}

% 表5 の傾向は図3 にも可視化されており、Modulo 精度と Magnitude 精度のスケーリング幅の差異が
% 折れ線グラフで対比的に示される。
The trends in Table~\ref{tab:scale_trend} are visualised in Fig.~\ref{fig:scaling}, which
contrasts the differential scaling of Modulo and Magnitude accuracy across model sizes.

\begin{figure}[t]
  \centering
  \includegraphics[width=\linewidth]{figures/fig3_scaling}
  \caption{Scaling behaviour across model sizes (Small / Middle / Large).
           Left: Magnitude accuracy (Mag Acc) remains high for both IntSeqBERT and Ablation,
           with only $+1.1$\,pt improvement from Small to Large.
           Centre: MMA (Mean Modulo Accuracy) for IntSeqBERT improves from 40.3\% (Small) to
           50.4\% (Large, $+10.1$\,pt), consistently exceeding Ablation at all scales.
           Right: Solver Top-1 accuracy for IntSeqBERT improves from 14.1\% to 19.1\%
           ($+5.0$\,pt), with a growing gap over Vanilla (${\approx}2.4$\%, flat).}
  \label{fig:scaling}
\end{figure}

\subsection{Limitations}
\label{sec:limitations}

% 大きな整数の予測困難。|x| ≥ 10^{20}（Huge・Astronomical バケット）では Solver 精度がほぼゼロ。
% ボトルネックは不十分な Modulo 予測精度にある。CRT による復元には複数の法に対する正確な剰余が必要
% であり、MMA が 50% 程度であっても剰余の誤りが頻発し、再構成が失敗する。
\textbf{Difficulty with large integers.} Solver accuracy collapses to near zero for
$|x| \geq 10^{20}$ (Huge and Astronomical buckets). The bottleneck is insufficient Modulo
prediction accuracy: CRT-based reconstruction requires \emph{exact} residues for multiple
moduli simultaneously, and with MMA at approximately 50\%, residue errors are frequent enough
to cause reconstruction failure even in the CRT mode (Top-1 = 0.09\%).

% 有効候補率の低下。IntSeqBERT の Solver 呼び出しの約 13% で有効候補が返らない（mode = none）。
% 主に CRT 失敗が原因。残差制約の緩和や近似 CRT の利用など Solver のフォールバック戦略の改善で
% 一部を回収できる可能性がある。
\textbf{Low valid-candidate rate.} Approximately 13\% of IntSeqBERT Solver calls return no
valid candidate (\texttt{none} mode), primarily due to CRT failures. Recovery may be possible
by relaxing residue constraints or adopting approximate CRT in Solver fallback strategies.

% データセットバイアス。OEIS の収録数列は nonn（非負・非全零）タグを持つものが多数を占める。
% 符号精度 98.54% という高い値の一部はこのバイアスに起因し、負の整数を多く含む数列に対しての
% 評価では下振れる可能性がある。
% 値が極めて大きい数列（天文学的バケット）は相対的に少数であり、そのバケットでのメトリクスは
% 少数サンプルで算出されている点にも留意が必要である。
\textbf{Dataset bias.} OEIS sequences are predominantly non-negative
(tagged \texttt{nonn}\hspace{0pt}). Our training and test data inherit this
distribution, so the high sign accuracy (98.54\%) may partially reflect this
bias and could degrade on sequences with many negative terms.
Additionally, Astronomical-bucket metrics are based on a small number of samples and should
be interpreted with caution.
